%% ==================================================================
%% From the research note "Where Does Agency Live?" --- the thesis,
%% with the reconciling material from "Not Every Cut Is an Aperture".
%% ==================================================================
\chapter{Where Does Agency Live?}
\label{ch:agy}

\noindent
Agency and determinism is an old argument, and one reason it does not move is that
both sides help themselves to the boundary of the agent as though it were given. Fix
the boundary and the question becomes tractable in an uninteresting way; leave it
floating and the question is unanswerable. The previous chapter supplies something
better than a stipulation. If choice is a resource and it is consumed at apertures,
then there is a non-arbitrary criterion for where to draw the line: draw it around the
regions that carry no observable choice.

That criterion has consequences, and this chapter is mostly a matter of following them.
The first is that the regions in question turn out to be small. If selection, predation
and competition are non-confluent --- and they are, almost by definition, since which
variant fixes and which lineage splits is exactly an outcome where the choice matters
--- then the determinate islands are local and the sea of races is dense. Apertures are
everywhere, and ``where is the agent?'' has many local answers rather than one global
one.

The second consequence is the one i did not expect. Decorating each cut with the
strength of choice it demands turns the factorisation into a \emph{filtration}: not one
partition into islands and seas but a family of them, one per level of choice strength,
coarsening as the level rises. An observer endowed with strength $C$ sees the world
island at $C$. Which means the stratification is not a tower imposed on the world from
outside; it is the level-sets of a single decorated term, and different observers are
different thresholds on the same object. This is also, and not by coincidence, what
consciousness turns out to be in the chapters that follow --- not a threshold an agent
crosses but a point it occupies, with the ordinal tower of \S\ref{sec:fibration} the
chain-shaped shadow of the lattice that point really lives in.

The chapter ends on a question rather than a result, and on a worry. The question is
whether the graded factorisation survives reduction --- whether an island can crack
while it runs, which would be the formal content of one agent perturbing another. The
worry is that if choice in nondeterministic interaction is how hypercomputation would
show up physically, it would be very hard to see. The closing section argues that hard
to see is not the same as invisible, and takes Bell as the precedent.

\bigskip
\section{Choice and the drawing of boundaries}\label{ndo:sec:boundaries}

We can now turn the instrument on agency. If nondeterminism is localised at cuts
and is a resource, then the natural place to look for an agent's boundary is where
its choices are live --- which is to say, at the apertures of
Definition~\ref{ndo:def:aperture}, governed by the principle that not every cut is
one.


\noindent
This is already a fresh answer to the old question. The apertures are exactly the
races; the determinate traffic between them is not where agency is decided, however
much computation flows through it. And it locates interference precisely: a
high-capacity agent reaches into a low-capacity one not everywhere they touch but
only at the cuts where the low-capacity one has a live, observable choice that the
high-capacity one can see how to resolve. The asymmetry of capacity --- one agent
several halting-oracles above another --- becomes operative only through an aperture.

Two corrections to a tempting oversimplification, both of which we will need. First,
the aperture is a property of the \emph{cut}, not of either agent; nondeterminism is
not preserved by $\Par$, so a determinate whole can have wildly underdetermined
internal cuts, and conversely. ``An agent's nondeterminism budget is its interface''
is therefore not quite right: the budget belongs to the cut, and which cut one calls
the boundary is a modelling decision. Second, because of this, a seam where choice is
genuinely open may sit far from the boundary an observer would draw, and its effects
may be felt at cuts the observer counts as internal. We turn to a setting where this
plays out concretely.


\section{The biological axis: Smith redraws the boundary}\label{ndo:sec:smith}

Smith and Morowitz read terrestrial life as a planetary process and ask after its
interface to the rest of the universe \cite{SmithMorowitz}. On their account the agents that
matter at that interface are the ones that consume energy from \emph{outside} the
biosphere --- sunlight, the chemical disequilibria of hydrothermal vents --- while a
vast interior of agents consume \emph{one another}, trading energy already captured.
The energy-eaters are the actual boundary between (terrestrial) life and the cosmos,
and Smith presses the point we want: that we may have to be
\emph{individual-organism-blind} and treat the entire biosphere as the agent. The
naive ontology places the agency boundary at the organism --- this bee, that bird ---
and Smith argues that on a thermodynamic criterion the boundary belongs somewhere else
entirely, around the whole. For him the biosphere is the macro-component, unified by
its single interface to the extra-biospheric energy source.

This boundary-redrawing is exactly the move we need, and it is the only thing we take
from the analogy. We do \emph{not} take Smith's interior to be determinate. It plainly
is not. Schematically,
\[
\textsf{Biosphere} \;\rcong\; \textsf{EnergyEaters} \Par \textsf{AgentEaters},
\quad
\textsf{Universe} \;\rcong\; \textsf{Biosphere} \Par \textsf{RestOfUniverse},
\]
the seam to the cosmos being $\textsf{Biosphere} \Par \textsf{RestOfUniverse}$ and
\textsf{EnergyEaters} the sub-term whose business is that cut --- where outside energy,
untyped to any particular molecule (the sun does not address its photons to a chosen
chlorophyll), enters as an underdetermined produce, a genuine aperture. But
\textsf{AgentEaters} is \emph{full} of meaningful races, exactly the ones
\S\ref{ndo:sec:interaction} put on the table: predation contended (two predators, one
prey), competition for scarce resource, and, most decisively, \emph{natural selection
and speciation}, which are non-confluent almost by definition --- which variant fixes,
which lineage splits, is precisely an outcome where the choice matters and the
histories do not reconverge. Smith's biosphere-agent is therefore \emph{not} a
macro-component in our sense. Its interior runs the deepest races biology has.

Those races are not new to this book. The chapters on ecology took the same
distinction --- organisms that eat only external energy against organisms that eat one
another --- and derived it as an \emph{access profile} rather than a kind of substance
(Definition~\ref{sci:def:source} and Proposition~\ref{sci:prop:twosciences} of
\chSci), with contended predation as a race on a metabolic channel and selection as
the mechanism by which a lineage reaches what no individual can
(\S\ref{sci:S12} of \chSci). That two developments starting from different premises
--- one from conserved tokens, one from confluence --- arrive at the same seam is the
kind of agreement worth noticing, and it is not a coincidence: a linear send that can
be received by either of two harvesters \emph{is} a contended channel, and the
foraging race is the aperture.

So the relationship between Smith's factorisation and ours is not that we fill in his
interior; it is that we propose a \emph{different criterion} for the same kind of move.

\begin{remark}[Two boundary-redrawings; ours strictly finer]\label{ndo:rmk:refine}
Smith redraws the naive per-organism boundary \emph{outward}, to one biosphere-sized
agent, on a thermodynamic criterion: what interfaces with the external energy source.
We redraw it by a \emph{different} criterion --- confluence up to resource bisimilarity
(\S\ref{ndo:sec:macro}) --- which lands the boundary neither at the organism nor at the
whole, but around whatever subsystems carry no observable choice. This is strictly more
refined: it asks not ``what touches the outside?'' but ``where is there no live race?''
The two need not nest. Smith's macro-component contains selection, so it is not one of
ours; our macro-components, as \S\ref{ndo:sec:macro} argued, are small and local, so the
biosphere is not one of theirs. Smith's whole and our islands are different objects,
unified by different things --- thermodynamics versus determinacy --- and the honest
statement is that the factorisations \emph{cross} rather than refine one within the
other. What we credit Smith with, and it is the load-bearing credit, is the
demonstration that the naive placement of the agency boundary is \emph{negotiable}:
``agent'' need not mean ``organism.'' Once that is conceded, the question is which
criterion should fix the boundary, and we answer it differently.
\end{remark}

\noindent
This still gives the seam-versus-boundary phenomenon of \S\ref{ndo:sec:boundaries} its
biological face, but more carefully than a single clean interface would. The energy cut
is an aperture; its bias propagates into a teeming interior that has plenty of its own
apertures; and the genuine computational question --- whether such a bias is, at any
given depth, transmitted through a determinate region or absorbed into a fresh race ---
is answered locally, region by region, not once for the whole. To answer it at all we
need the determinate regions named precisely, and they are smaller than Smith's whole.


\begin{remark}[An earlier reading, and why it does not survive]\label{ndo:rmk:earlier}
It is worth recording a reading this one displaces, because the displaced version is
the tempting one and its failure is instructive. Read quickly, Smith's factorisation
looks like a macro-component outright: energy-eaters at the seam, agent-eaters as a
rigid interior, the whole thing a determinate island with a single aperture onto the
sun. The interior even looks the part, being spectacularly constraint-sensitive ---
one substrate per active site, this bee to that flower. But conformance constraints
are not confluence. An enzyme's specificity says which pairings are \emph{admissible};
it says nothing about which of several admissible pairings occurs, and it is the latter
that Definition~\ref{ndo:def:macro} asks about. Two predators and one prey is a fully
type-correct configuration and a live race. The interior is constrained and
non-confluent at once, which is precisely the combination the rigidity criterion is
built to detect and the linearity criterion would miss.
\end{remark}

\section{Decorating the cuts: the factorisation, graded by choice}\label{ndo:sec:graded}

The companion development \cite{choicesched} assigns to the nondeterminism at a cut a
\emph{choice type}: the strength of choice principle --- a degree in the Weihrauch lattice
\cite{weihrauch}, equivalently an altitude in the stalk of \S\ref{sec:fibration} --- required to
resolve that cut's race, fairly and correctly, taken modulo resource bisimilarity so that
branching which reconverges costs nothing. Write $\chi(c)$ for the choice type of a cut $c$, and
$\bot$ for the bottom of the lattice: no choice, computable resolution. Then the two readings of
\S\ref{ndo:sec:macro} collapse into a single labelling,
\[
\chi(c) = \bot \quad\Longleftrightarrow\quad c \text{ is not an aperture,}
\]
a cut being interior to a determinate region exactly when its choice type is trivial --- whether
because the match is unique or because all matches are $\rbisim$-equal, proof-irrelevant. The
macro-component factorisation of \S\ref{ndo:sec:macro} is read straight off the decoration: the islands
are the maximal subterms whose internal cuts all carry $\bot$, the seas the cuts that carry more.

Once the cuts are labelled, the binary distinction --- $\bot$ against more-than-$\bot$ --- is
revealed as the bottom case of a graded family. Fix any level $C$ in the lattice.

\begin{definition}[$C$-island]\label{ndo:def:cisland}
A \emph{$C$-island} is a maximal subterm all of whose internal cuts $c$ satisfy $\chi(c) \le C$:
every internal race is resolvable within choice resources of strength $C$. The cuts with
$\chi(c) \not\le C$ are the \emph{$C$-seas}; in the chain case these are the cuts of strictly
greater strength $C' > C$, demanding computation higher in the hypercomputation tower than $C$
reaches.
\end{definition}

\noindent
A $C$-island is determinate \emph{to an observer endowed with choice strength $C$}: such an
observer resolves everything inside and sees no open choice, and must climb only at the seas, where
the type exceeds $C$. This is the exact form of a relativity the inquiry has circled from the start
--- ``determinate relative to a higher resource'' --- now indexed by a definite level, the
observer's own altitude. Determinacy is not absolute; it is relative to the choice power brought to
bear, and the threshold $C$ is that power named.

\begin{proposition}[The factorisation is a filtration]\label{ndo:prop:filtration}
If $C \le C'$ then every $C$-island is contained in a $C'$-island: raising the threshold coarsens
the partition, merging islands across any sea whose strength lies between. At $C = \bot$ the
partition is the confluent-islands / non-confluent-seas factorisation of \S\ref{ndo:sec:macro}; at
$C = \top$ the whole term is a single island. A term therefore carries not one factorisation but a
\emph{filtration by choice strength} --- a family of factorisations, monotone in the threshold,
indexed by the lattice.
\end{proposition}

This is where the flat ontology and the tower meet, and it is worth saying exactly how. In
\S\ref{sec:fibration} the stratification was \emph{external}: a tower indexed over a category of models,
a tower projected onto a term from outside. Here it is \emph{internal}. The filtration
$\{C\text{-islands}\}_{C}$ is the term's own stratification, read off the decoration $\chi$ of its
cuts, not imposed on it. The tower is the decoration. And the flatness is precisely that there is
one term carrying one labelling, of which the strata are the level-sets: different observers are
different thresholds, each seeing a different islanding of the same decorated term. ``The tower is a
cleavage of a flat jumble'' (\S\ref{ndo:sec:flat}) sharpens into the statement that the cleavage
\emph{is} $\chi$, and the jumble is the single $\chi$-decorated term beneath all of its thresholds.

\begin{remark}[Do Smith's whole and our islands nest after all?]\label{ndo:rmk:nest}
At $C = \bot$ Smith's biosphere is not one of our islands: selection and predation are seas
(\S\ref{ndo:sec:smith}). The filtration reframes the question. There may be a level $C^{*}$ --- the
strength at which selection-type races become resolvable --- at which the biosphere \emph{is} a
$C^{*}$-island, so that Smith's thermodynamic whole nests inside our graded factorisation at
altitude $C^{*}$ even though it does not at the bottom. Whether such a uniform $C^{*}$ exists ---
whether the biosphere's races share a single choice degree or scatter across the lattice --- we do
not know. But the filtration turns ``do they nest?'' from a verdict into a measurement: \emph{at
what level, if any?} That is the better question, and it is the decoration that poses it.
\end{remark}

\noindent
There is a second reading of the filtration worth recording, because it connects this
part to the one before it. Proposition~\ref{sci:prop:confinement} of \chSci says an
individual learner revising in small steps can never leave the ball it started in, and
Corollary~\ref{sci:cor:poverty} says the poorer it is the smaller that ball. Read
through the decoration, confinement is a statement about \emph{threshold}: a learner
whose budget affords only choice strength $C$ sees a $C$-islanding of its world and
cannot even formulate a hypothesis that would require distinguishing what lies inside
one of its islands. Poverty is confining because poverty is a low coordinate. And the
escape found there --- that the unit of learning must move from the individual to the
lineage to the composite --- reads here as raising $C$ by acquiring a stronger
scheduler, which is why composition and merger were the operations that worked.

Two honesties temper the construction. First, the decoration is well defined on a static term only
insofar as subject reduction holds (\cite{choicesched}; \S\ref{ndo:sec:crux}): if scheduling can climb
the stalk, $\chi(c)$ rises as the term runs, an interior cut can cross above $C$ and split its
island, and the filtration is not invariant under reduction. \emph{Perturbation, in this language,
is an island cracking as it runs} --- the $C$-factorisation refining dynamically because a boundary
decision manufactured strength inside. The crux of \S\ref{ndo:sec:crux} is therefore exactly the
question whether the graded factorisation is a reduction invariant. Second, $\chi$ is a semantic
invariant, not in general computable --- Weihrauch degree is undecidable at large. But it is
approximable from above by \emph{types}: linear internal typing forces $\chi = \bot$, session typing
bounds it, finite contention keeps it constructive. The term assignment of the companion note is the
effective approximation of the decoration, which is why a factorisation defined semantically can
nonetheless be certified syntactically.


\section{Stratification versus the jumble}\label{ndo:sec:jumble}

The natural ontology suggested by ``choice imports computational power'' is a
\emph{tower of oracles}: agents stratified by degree, the stronger able to reach
down into the weaker. We want to register a different intuition and then say
exactly what saves it. A universe made of computational phenomena may be not a
neat tower but a \emph{jumble} --- agents of all stripes and capabilities
interacting, as we see in the physical world, with the stratification only a
framework that aids understanding, not the structure of the world. If that is the
case, agents with strong oracles could interfere with the ``choices'' of agents
less computationally endowed, but they do so as participants in a flat field, not
as occupants of a higher floor.

\paragraph{Both pictures at once: a (bi)fibration.}
The two readings need not compete. Let a total category $\mathcal{E}$ be the
jumble --- every agent of every degree, with the cut-mediated maps between them ---
and let a projection to the degrees (Turing, Weihrauch, or an ordinal index) be
the stratification. Reindexing along $\dgr' \le \dgr$ --- the cartesian, fibration
direction --- takes a degree-$\dgr$ computation to its degree-$\dgr'$ shadow, the
view a weaker observer keeps after forgetting the oracle. That is the
``stratification as aid to understanding'' move made precise: \emph{the tower is a
cleavage, a chosen gauge for describing a frame-independent jumble}. Interference
is the other direction, and for it one wants the reindexing functors to carry
Lawvere adjoints $\exists_f \dashv f^{*} \dashv \forall_f$. We record, as a
conjecture rather than a theorem, the fit that motivates the formalisation:

\begin{conjecture}[Angelic/demonic as adjoints]\label{ndo:conj:adjoint}
Angelic resolution (the cooperative scheduler that \emph{finds} a good pairing)
and demonic resolution (the adversarial scheduler that must survive \emph{every}
pairing) are, respectively, the left and right Lawvere adjoints
$\exists_f \dashv f^{*} \dashv \forall_f$ to reindexing along the
degree/prefix order. If so, the scheduler ladder --- canonical, demonic-finite,
angelic, fair-merge, relativised, transfinite --- is the adjoint structure of one
fibration rather than a list of cases.
\end{conjecture}

\noindent
The load-bearing check is whether the fibres carry the limits and colimits to
support both adjoints; that is where a formalisation should begin.

\paragraph{What keeps the jumble from collapsing.}
A flat ontology of \emph{unequal} agents is only coherent if a single strong agent
cannot simply dominate. Two design decisions of the source vocabulary are exactly
the consistency conditions. \emph{No implicit interaction} \cite{paths}: the store
never reacts with itself; a reaction is triggered only by a cut that actually
co-locates opposite polarities. Drop it, and a single halting-oracle agent on a
universal channel homogenises the field to the top degree --- the jumble dies.
\emph{Locality of the cut}: interference can be written only where a channel is
shared, and, in the path reading, only into the subspace below a held prefix.
Position in the trie and degree \emph{together} set the reach of interference. Add
a cost discipline (the phlogiston of the source programme) and reach becomes finite
per unit resource. Locality plus cost is what keeps a multi-degree universe a
jumble rather than a hierarchy that flattens upward. We state this as a slogan,
because it is the most decision-relevant thing the vocabulary offers: \emph{no
implicit interaction and the risk ledger are not bookkeeping; they are the
consistency conditions for a flat ontology of unequal computational agents.}


\section{The flat ontology}\label{ndo:sec:flat}

The tower of \S\ref{sec:fibration} is a map, and we should not mistake it for the
territory. It is a way to \emph{understand} computational phenomena --- to sort agents
by capacity and read interference as reaching down a tower. But the actual universe of
computational phenomena may be far less tidy, and biology is the reason to suspect so.
In a single living cell, fast molecular dynamics --- vibrational and electronic
transitions down to the femtosecond --- coexist with transcription and translation
that unfold over seconds to minutes: many orders of magnitude apart, intimately
coupled, in one small volume. The cell does not run its fast and slow processes in
separate strata that communicate through a clean interface; it juxtaposes them, and the
juxtaposition is the point. The physical universe at large is the same --- scales and
scopes wildly interleaved rather than cleanly layered.

We propose to take computation the same way. In the \emph{flat ontology}, the universe
of computational phenomena is the total space of the fibration \emph{without privileging
the projection} --- a jumble of agents of wildly different computational capacity,
several halting-oracles up beside fully deterministic $\lambda$-like agents,
interacting directly at cuts. The tower is not the structure of the world; it is a
cleavage, a chosen gauge for describing a frame-independent jumble, and the
``stratification'' is an artefact of which morphism we chose to project along --- or, read
internally, of which threshold in the filtration of \S\ref{ndo:sec:graded} we choose to view the term
at. The
islands-in-a-sea-of-races picture of \S\ref{ndo:sec:macro} is the same claim seen from the
agency side: determinate regions are local and sparse, apertures pervade, and there is no
clean global stratum at which all the choice has been pushed to one rim. Read this in
both directions, as promised. Computation, taught by biology, stops expecting a neat
hierarchy and starts expecting juxtaposition. Biology, taught by computation, gains a
vocabulary for what the juxtaposition costs and where it can bite: a high-capacity agent
can influence a low-capacity one, but --- Principle~\ref{ndo:prin:main} --- only at an
aperture, only where the lower agent has a live choice the higher one can resolve.

What keeps such a flat universe coherent, rather than collapsing to the top capacity?
Two disciplines, both already in the MPC vocabulary. \emph{No implicit interaction}: the
store never reacts with itself; a reaction is triggered only by a cut that actually
co-locates opposite polarities \cite{paths}. Drop it and a single oracle on a universal
channel homogenises the field upward; keep it and interference must be \emph{cut-triggered}.
\emph{Locality and cost}: a strong agent can write only where it shares a channel ---
and, in the path reading, only into the subspace below a prefix it holds --- and a cost
discipline makes its reach finite per unit resource. Locality plus cost is what lets
agents of disparate capacity share a world without the strongest simply absorbing the
rest. These are not bookkeeping; they are the consistency conditions for a flat ontology
of unequal agents --- and they are exactly the conditions under which a small confluent
island can persist amid the surrounding races without being dissolved by the strongest
agent that happens to share one of its channels.


\section{The hinge: is determinism stable under interference?}\label{ndo:sec:crux}

One question decides whether the picture is static or dynamic, and it is the question a
formalisation should resolve first. Everything above treats ``determinate interior,
nondeterministic boundary'' as a property a system \emph{has}. Whether it \emph{keeps}
the property when a boundary aperture fires and the decision propagates inward is not
obvious --- because confluence up to resource bisimilarity, unlike linearity, is not
compositional, and because a deep interaction in the path semantics manufactures a fresh,
deeper interaction whose confluence is a new obligation \cite{paths}.

\begin{question}[Stability of rigidity]\label{ndo:q:crux}
Is confluence up to resource bisimilarity preserved as a boundary decision cascades into
the interior? Equivalently: when a macro-component's aperture fires, does its interior
stay rigid?
\end{question}

\noindent
The two answers are both meaningful, and the choice between them is what biology and
computation are jointly asking. If rigidity is stable, the factorisation is robust: a
boundary race fires, the decision descends as a \emph{forced} cascade through a confluent
island that has no opinion of its own, and ``determinate interior / nondeterministic
boundary'' is a genuine structure theorem holding along the dynamics --- a buffered
developmental cascade absorbing an upstream choice into one canalised outcome. If rigidity
is \emph{not} stable, then boundary interference can crack open a previously determinate
island, turning a confluent agent into a branching one by descending into it. That is not
merely a defect;
it is plausibly the formal content of \emph{perturbation} --- a high-capacity agent not
only steering a low-capacity one within its existing freedom but \emph{manufacturing} new
freedom in it and then steering that. In the ontological register: the difference between a
universe whose agents have fixed capacities and one in which capacity itself is something
agents can do to one another. We do not know which holds, and the same confluence lemma
settles it on both axes at once.


\subsection{The crux, in proof-theoretic dress}

The term assignment of \chChs\ gives this question a second face, and the two faces
are the same lemma.

\begin{question}[The hinge, as subject reduction]\label{ndo:q:sr}
Does the term assignment satisfy subject reduction across the cascade in which one
interaction manufactures a deeper one? Equivalently: can the resolution of a boundary
race cause an interior to acquire a choice type it did not have, climbing the stalk as
it reduces?
\end{question}

\noindent
If subject reduction holds, the graded factorisation is stable. Determinate interiors
stay determinate under scheduling, and a scheduler's imported power is fixed by its
type once and for all. If it fails, the failure \emph{is} perturbation: a high-capacity
scheduler, reducing against a low-capacity system, manufactures new choice in it --- a
well-typed term reducing to an ill-typed one, which is the type-theoretic signature of
one agent changing another's capacity. That a single metatheorem settles both the
agency question and the metatheory of the correspondence is the best evidence we have
that the correspondence is the right one. \chChs\ develops the correspondence, and
what it owes.

\section{Open problems}\label{ndo:sec:open}

We collect the load-bearing uncertainties, each a structural claim one can try to settle.
\emph{(i)} Make precise the family-of-pairings whose section is the schedule and prove the
$\mathrm{AC}$-grading of \S\ref{ndo:sec:choice}, identifying which schedulers are confluent up
to $\rbisim$ (constructive, no imported oracle) and which import altitude in the stalk.
\emph{(ii)} Exhibit the fibration $p\colon\Hyp\to\mathbf{GSLT}$ rigorously and test whether angelic
and demonic resolution are the Lawvere adjoints $\exists \dashv (-)^{*} \dashv \forall$ to
reindexing --- the decisive sub-question being whether the GSLT fibres carry the limits and
colimits both adjoints require. \emph{(iii)} Show ``maximal rigid subsystem''
(Definition~\ref{ndo:def:macro}) is well defined and that the three frontiers of
Proposition~\ref{ndo:prop:frontier} coincide, and establish that a confluent interior is
resource-degree-transparent. \emph{(iv)} Resolve Question~\ref{ndo:q:crux} --- the stability of
rigidity under descent --- proving the structure theorem if it holds, or characterising
perturbation as a property of the sub-system if it fails. \emph{(v)} Settle the scale
question of \S\ref{ndo:sec:macro}: is any biosphere-scale structure confluent up to $\rbisim$,
or is confluence strictly small-scale with all large structure race-dominated? If the
latter, confirm that Smith's thermodynamic macro-component and our determinacy
macro-components \emph{cross} rather than nest, and characterise each as the carrier of a
different invariant --- and, via the filtration of \S\ref{ndo:sec:graded}, determine whether they
nonetheless nest at some higher level $C^{*}$ (Remark~\ref{ndo:rmk:nest}). \emph{(vi)} Test the
candidate that some local confluent island sits at ecosystem scale against a concrete ecological
model, and locate the apertures that bound it. \emph{(vii)} Develop the choice-type decoration
$\chi$ and the filtration by choice strength of \S\ref{ndo:sec:graded}: prove the monotonicity of
Proposition~\ref{ndo:prop:filtration}, establish that the $C$-factorisation is a reduction invariant
\emph{iff} the term assignment of \cite{choicesched} satisfies subject reduction (so that
``perturbation'' is exactly an island cracking under reduction), and pin down the syntactic
approximations (linear and session types as upper bounds on $\chi$) that make the semantic
factorisation effectively certifiable.


\section{Coda: hard to see, and the precedent of Bell}\label{ndo:sec:coda}

If this picture is right --- if, in the physical world, the surplus power of a high-capacity
agent shows up as nothing more dramatic than the way a low-capacity agent's nondeterministic
interactions get resolved --- then hypercomputation in nature would be extraordinarily hard to
see, and the difficulty would be principled rather than incidental. A confluent region hides
its choices by construction (\S\ref{ndo:sec:macro}): the resolution happens but reconverges,
leaving no trace. And even at an aperture, a single resolution looks like nothing --- one
outcome among those the situation allowed, indistinguishable, taken alone, from a coin flip.
Choice-as-resource emits no per-event signal. Whatever evidence there is cannot live in any
single outcome; it must live in the \emph{correlations across} outcomes. The signature of a
resolution stronger than a local choice function is that separated apertures come out
correlated in ways no local, independently-resolving assignment could produce. That is why the
resource is hard to test --- and it tells us exactly where to look.

It is, almost word for word, the lesson of Bell. Bell's theorem shows that no theory in which
each measurement outcome is fixed by local, pre-assigned values --- no \emph{local} choice
function for the outcomes --- can reproduce the correlations quantum mechanics predicts for
separated measurements on an entangled pair \cite{bell}; and the inequalities that mark the
boundary have been violated in the laboratory, culminating in loophole-free experiments
\cite{hensen}. Read in our vocabulary, a quantum measurement is an aperture: a cut whose
outcome the local term does not fix. Bell is then a statement about \emph{correlated choice} ---
that the joint resolution of two separated apertures carries a correlation no local scheduler,
no locally-supplied choice function, can account for. The per-event view sees only randomness
at each wing; the structure is in the correlation, precisely as the testability argument above
predicts. And the correlations have a developed mathematics: the device-independent resource
theory of nonlocality treats these correlations as a graded resource \cite{brunner}, and the
sheaf-theoretic account of contextuality \cite{abramsky} --- logic-and-computer-science
apparatus, already in our bibliography --- characterises exactly the obstruction to a global,
locally-consistent assignment. The bridge between the choice mathematics and the physics already
has a pier on the far bank.

We are not claiming that Bell violations exhibit hypercomputation, or that nonlocal correlations
compute anything uncomputable: they do not, and the resource that nonlocality supplies is not
Turing degree. The precedent is methodological, and in spirit it is decisive. It establishes
that \emph{correlated choice is a real, mathematically characterised, experimentally confirmed
feature of the physical world} --- that a phenomenon which is invisible per-event, visible only
in correlation, and naturally described as the nonlocal resolution of an otherwise-underdetermined
interaction, is not a fantasy of testability but the best-confirmed strange fact in physics.
A research posture that treats choice as a physical resource, and looks for its signature in
correlations across apertures rather than in single events, has therefore already succeeded once,
spectacularly. That is reason to think it could succeed again, for resources further up the
hierarchy than nonlocality reaches.

This is the note we want to leave the reader with, and it points in both directions, as the whole
essay has tried to. Bringing the mathematics of choice --- the axiom of choice and its gradations,
the Weihrauch lattice, the resource theories --- into contact with a physics informed by computation
and biology may be fertile for each. The mathematics gains a fresh set of intuitions and motivations:
choice stops being only a question in the foundations of sets and becomes a calculus of physical and
computational resource, with the strength of a choice principle reading as the altitude it buys in a
hierarchy of capacity, and with empirical stakes. The physics gains a new set of tools: the category
of models and its fibration of hypercomputation, choice-as-resource as the vertical coordinate, the
aperture and the macro-component as a vocabulary for where in a system agency and capacity actually
sit, and a flat ontology in which entities of wildly unequal computational power --- a halting oracle
and a $\lambda$-term --- can be imagined to meet, as so much else in nature meets, at a single cut.

\paragraph{Acknowledgements.}
This essay draws on the RSpace denotation of the rho calculus \cite{qcs,paths} and the GSLT/OSLF
programme \cite{OSLF}, and on Smith and Morowitz's reading of the biosphere \cite{SmithMorowitz}; its
conceptual claims are deliberately routed back to structural obligations where they can be settled.
